\documentclass{article}
\usepackage[T1]{fontenc}
\usepackage{lmodern}
\usepackage{graphicx}
\usepackage{amsmath,amssymb}
\usepackage[a4paper,margin=2cm]{geometry}
\usepackage[absolute,overlay]{textpos}
\setlength{\TPHorizModule}{1mm}
\setlength{\TPVertModule}{1mm}
\usepackage[indonesian]{babel}
\usepackage{hyperref}
\setlength{\emergencystretch}{2em}
% unit: Halaman 24

\begin{document}

% === Halaman 24 (llm) ===

\section*{Kurva Eliptik pada Galois Field}

\begin{itemize}
    \item Operasi kurva eliptik yang dibahas sebelum ini didefinisikan pada bilangan riil.
    \item Operasi pada bilangan riil tidak akurat karena mengandung pembulatan
    \item Pada sisi lain, kriptografi dioperasikan pada ranah bilangan integer.
    \item Agar kurva eliptik dapat dipakai di dalam kriptografi, maka kurva eliptik didefinisikan pada medan berhingga atau Galois Field, yaitu $\text{GF}(p)$ dan $\text{GF}(2^m)$.
    \item Yang dibahas dalam kuliah ini hanya kurva eliptik pada $\text{GF}(p)$
\end{itemize}

\end{document}
